% =====================================================================
% §2 Setup — task, data, model, deployment
% =====================================================================

\section{Setup}
\label{sec:setup}

\paragraph{Task.}
At forecast issuance time $T$, predict hourly per-zone MWh demand for
the 8 ISO-NE load zones over the next 24~hours. With $S=24$ history
steps and 44-d calendar one-hots $C_t$, the model computes
\begin{equation}
\hat Y_{t+1:t+24} = f_\theta(X_{t-S:t+24} \in \mathbb{R}^{(S+24)\times 7\times 450 \times 449},\;
   Y_{t-S:t} \in \mathbb{R}^{S\times 8},\;
   C_{t-S:t+24} \in \mathbb{R}^{(S+24)\times 44}),
\label{eq:task}
\end{equation}
where $X$ are 7-channel HRRR weather rasters
\citep{benjamin2016hrrr,noaa2024hrrr}. The metric is MAPE in physical
MWh space averaged over the 24-h horizon and 8 zones. All numeric
claims in the main paper come with bootstrap 95~\% confidence intervals
over (forecast-day, zone) pairs (procedure in
Appendix~\ref{app:bootstrap}).

\paragraph{Data and splits.}
ISO-NE per-zone demand is from the public Energy / Load / Demand
reports portal \citep{isone2024data}. HRRR weather is the f00 analysis
product on a $450 \times 449$ NE bbox, regridded from native Lambert
conformal via Delaunay barycentric interpolation
(Appendix~\ref{app:model}). At training time both the history and
future weather windows use f00 analyses; at deployment we substitute
HRRR f01..f24 forecasts for the future window, and quantify the
resulting OOD penalty in Appendix~\ref{app:forecast_shift}. Splits are
listed in Table~\ref{tab:splits}.

\begin{table}[h]
\centering\small
\caption{Train / val / test splits used in this paper.}
\label{tab:splits}
\begin{tabular}{@{}llr@{}}
\toprule
Split & Years & Approx. samples \\
\midrule
Train (baseline)               & 2019--2020              & $\sim 17\,$K \\
Validation                     & 2021                    & $\sim 8\,$K  \\
Test A (training-window self-eval) & 2022-12-30 to 12-31  & 2 forecast days \\
Test B (deployment W1)         & 2025-05-01 to 05-07     & 7 forecast days \\
Test B (deployment W2)         & 2026-04-28 to 05-04     & 7 forecast days \\
\bottomrule
\end{tabular}
\end{table}

\paragraph{Model.}
The trained baseline is a 1.75-M-parameter CNN-Transformer that fuses
HRRR weather, 24-h per-zone demand history, and 44-d calendar one-hots
in a single 4-layer Transformer encoder (architecture details in
Appendix~\ref{app:model}). For deployment we additionally blend with
\textbf{Chronos-Bolt-mini} \citep{ansari2024chronos} (Amazon, 21~M
params, Apache-2.0) used \emph{zero-shot} on 672~h of per-zone demand
history alone (no weather, no calendar). The two predictions are
combined per-zone:
\begin{equation}
\hat Y_{t+h, z} = \alpha_z \cdot \hat Y^{\text{baseline}}_{t+h, z}
+ (1 - \alpha_z) \cdot \hat Y^{\text{Chronos}}_{t+h, z},
\label{eq:ensemble}
\end{equation}
with fixed $\alpha_z$ tuned on a 14-day 2022 validation window.
$\alpha_z = 0$ for CT, SEMA, NEMA\_BOST (deployment drops the baseline
in those zones); $\alpha_z = 0.80$ for VT (baseline-dominant). On the
2022 self-eval, the ensemble achieves 4.21~\% MAPE vs the baseline's
5.24~\%.

\paragraph{Public deployment.}
We deploy three independent public artifacts (Figure~\ref{fig:3repo}):
a code repo holding training source and the Gradio app source, a data
repo whose only purpose is to host daily-refreshed JSON outputs, and
the HF Space \texttt{jeffliulab/predict-power} that serves user
requests. A GitHub Actions cron in the data repo runs the rolling
7-day backtest every 24~h using the same pipeline the Space serves at
request time. Splitting code and data lets daily backtest refreshes
avoid triggering Docker rebuilds of the Space and keeps the code repo's
git history clean. Reproducibility URLs and build commands are in
Appendix~\ref{app:repro}.

\begin{figure}[h]
\centering
% TikZ — three-repo deployment architecture (used in §8.1).
\begin{tikzpicture}[
  font=\small,
  every node/.style={align=center},
  repo/.style={
    rectangle, rounded corners=3pt,
    draw=NavyAccent, thick,
    fill=CardBG,
    text=InkColor,
    minimum width=4.4cm, minimum height=1.4cm,
    inner sep=4pt,
  },
  source/.style={
    rectangle, rounded corners=2pt,
    draw=SubInk, dashed,
    fill=white,
    text=SubInk,
    minimum width=4.0cm, minimum height=0.8cm,
    inner sep=3pt,
    font=\footnotesize,
  },
  arrow/.style={-Stealth, thick, draw=AccentBlue},
  thinarrow/.style={-Stealth, semithick, draw=SubInk},
  edgelabel/.style={
    font=\scriptsize\itshape, text=AccentBlue,
    inner sep=1pt,
    fill=white,
  },
]

% --- Public data sources (top row) ---
\node[source] (hrrr)    at (-4.6, 5.2)  {\textbf{NOAA HRRR}\\\texttt{s3://noaa-hrrr-bdp-pds/}};
\node[source] (isone)   at ( 0,    5.2) {\textbf{ISO-NE}\\\texttt{transform/csv/...}};
\node[source] (chronos) at ( 4.6,  5.2) {\textbf{HuggingFace Hub}\\\texttt{amazon/chronos-bolt-mini}};

% --- Three repos (middle row) ---
\node[repo] (code)  at (-4.8, 2.4) {\textbf{Code repo}\\\texttt{real-time-power-predict}\\\footnotesize source-of-truth};
\node[repo] (data)  at ( 0,    2.4) {\textbf{Data repo}\\\texttt{...-data}\\\footnotesize cron output, JSON/CSV};
\node[repo] (space) at ( 4.8,  2.4) {\textbf{HF Space}\\\texttt{predict-power}\\\footnotesize Gradio serving};

% --- User (bottom) ---
\node[draw=SubInk, thick, fill=white, rounded corners=2pt,
      minimum width=4.0cm, minimum height=0.8cm,
      font=\footnotesize] (user) at (4.8, -0.2)
      {\textbf{User browser}\\real-time forecast + 7-day backtest};

% --- Source -> data-repo arrows (top) ---
\draw[thinarrow] (hrrr.south)    -- node[edgelabel, sloped, pos=0.55] {Herbie} ($(data.north) + (-0.6, 0)$);
\draw[thinarrow] (isone.south)   -- node[edgelabel, pos=0.50]         {cookie-prime CSV} (data.north);
\draw[thinarrow] (chronos.south) -- node[edgelabel, sloped, pos=0.55] {from\_pretrained} ($(data.north) + (0.6, 0)$);

% --- Source -> Space arrows (used at request time) ---
\draw[thinarrow] (hrrr.south)    .. controls +(-1.0, -1.5) and +(0, 2.0) .. ($(space.north) + (-0.6, 0)$);
\draw[thinarrow] (isone.south)   .. controls +(0, -1.5)    and +(0, 2.0) .. (space.north);
\draw[thinarrow] (chronos.south) .. controls +(0.6, -1.5)  and +(0.6, 2.0) .. ($(space.north) + (0.6, 0)$);

% --- Code -> Data (cron checkout) ---
\draw[arrow] (code.east) -- node[edgelabel, pos=0.5] {GitHub Actions\\checkout main} (data.west);

% --- Code -> Space (HF -> Space sync) ---
\draw[arrow, dashed] (code.south) .. controls +(0, -1.6) and +(-1.0, -1.6) ..
   node[edgelabel, sloped, pos=0.45] {HF auto-sync on push} (space.south);

% --- Data -> Space (raw URL fetch on startup) ---
\draw[arrow] (data.east) -- node[edgelabel, pos=0.5] {raw HTTPS GET\\$\sim$250 KB at startup} (space.west);

% --- Space -> User ---
\draw[arrow] (space.south) -- node[edgelabel, pos=0.5] {Gradio HTTPS} (user.north);

% --- Section labels (left margin) ---
\node[font=\small\itshape, text=SubInk, anchor=east] at (-7.2, 5.2) {Public sources};
\node[font=\small\itshape, text=SubInk, anchor=east] at (-7.2, 2.4) {Three repos};
\node[font=\small\itshape, text=SubInk, anchor=east] at (-7.2, -0.2) {Consumer};

\end{tikzpicture}

\caption{Three-repo deployment architecture.
The cron in the data repo checks out the code repo, runs the rolling
backtest builder against the same public data sources the Space
consumes at request time, and commits the JSON+CSV results back. The
Space pulls via raw HTTPS, no auth.}
\label{fig:3repo}
\end{figure}
